\documentclass[11pt,a4paper,oneside]{book}
\usepackage[T1]{fontenc}
\usepackage[utf8]{inputenc}
\usepackage[italian]{babel}
\usepackage{lmodern}
\usepackage{microtype}
\usepackage{geometry}
\geometry{margin=2.7cm}
\usepackage{amsmath,amssymb,booktabs,longtable,array}
\usepackage{xcolor}
\usepackage{listings}
\usepackage{hyperref}
\usepackage[nameinlink,noabbrev]{cleveref}
\usepackage{enumitem}
\usepackage{fancyvrb}
\hypersetup{colorlinks=true,linkcolor=blue!45!black,citecolor=green!35!black,urlcolor=blue!55!black,pdftitle={Manuale Neuro-Symbolic Text-to-Gloss}}
\lstdefinestyle{t2g}{basicstyle=\ttfamily\small,breaklines=true,breakatwhitespace=false,columns=fullflexible,frame=single,showstringspaces=false,upquote=true}
\lstset{style=t2g}
\newcommand{\pathcode}[1]{\nolinkurl{#1}}
\newcommand{\method}[1]{\textsf{#1}}
\newcommand{\warningbox}[1]{\par\medskip\noindent\fcolorbox{red!60!black}{red!5}{\parbox{0.92\textwidth}{\textbf{Attenzione.} #1}}\par\medskip}
\title{Neuro-Symbolic Text-to-Gloss\\Manuale tecnico e operativo}
\author{Documentazione del progetto \texttt{neuro\_symbolic\_t2g}}
\date{Protocollo corrente Qwen2.5-0.5B}
\begin{document}
\frontmatter
\hypersetup{pageanchor=false}
\maketitle
\cleardoublepage
\hypersetup{pageanchor=true}
\chapter*{Come usare questo manuale}
Questo testo insegna a comprendere, eseguire, analizzare ed estendere il progetto English-to-ASL-gloss. Le descrizioni operative derivano dalla configurazione corrente sotto \pathcode{experiments/configs/qwen25-05b/}; non sostituiscono il controllo del file YAML risolto conservato con ogni run. Il manuale non dichiara risultati numerici della campagna corrente.
\tableofcontents
\mainmatter
\chapter{Panoramica e contratto scientifico}\label{chap:overview}
Il progetto traduce testo inglese in sequenze di glosse ASL. Una glossa è
un'etichetta discreta associata a un segno, non una rappresentazione completa di
spazio, simultaneità o componenti non manuali. Il verso è English-to-gloss sul
corpus ASLG-PC12 \cite{othman2012aslg,hf_aslg}; risultati gloss-to-English non
sono confrontabili.

\section{Cosa viene studiato}
Il backbone è \pathcode{Qwen/Qwen2.5-0.5B-Instruct} \cite{qwen25}. Si confrontano
base, SFT, GRPO e SFT seguito da GRPO, con LoRA/QLoRA, retrieval TF--IDF e
decodifica lessicalmente vincolata. ``Neuro-simbolico'' indica qui l'interazione
fra distribuzione neurale e strutture discrete; non implica una grammatica ASL
completa.

\section{Tassonomia da non violare}
\begin{center}\begin{tabular}{p{.20\textwidth}p{.68\textwidth}}\toprule
Asse & Significato\\\midrule
Metodo & \method{base}, \method{sft}, \method{grpo}, \method{sft-grpo}: sistema appreso e inizializzazione.\\
Prompt & \method{zero-shot} o \method{few-shot}/retrieval: solo condizionamento in ingresso.\\
Variante & \method{pda}, \method{hot}, \method{reward-*}: ablazione controllata.\\
Lifecycle & \method{train}, \method{baseline}, \method{ablation}, \method{probe}: tipo di job/artifact.\\\bottomrule
\end{tabular}\end{center}
Una modalità di prompt non è un metodo. Una baseline zero-shot è il metodo base
non addestrato con quel prompt. Trie è il percorso primario; PDA è un'ablazione
manuale. I probe congelati non sono obiettivi addestrabili.

\section{Come leggere e riprodurre}
I capitoli 3--8 spiegano dati e metodi, 9--11 protocollo e diagnostica, 12--15
codice e operazioni. Prima di una conclusione verificare config risolta, hash,
split, checkpoint, modalità prompt, decoding e versione metrica. Questo manuale
non inventa risultati correnti. \pathcode{docs/SOURCES.md} è il registro completo
delle fonti; la bibliografia finale contiene quelle citate nel testo.

\chapter{Fondamenti}\label{chap:background}
\section{Modello autoregressivo}
Dato un prompt $x$ e una glossa tokenizzata $y$, il modello fattorizza
$p_\theta(y\mid x)=\prod_t p_\theta(y_t\mid x,y_{<t})$. Il tokenizer opera su
subtoken, mentre metriche e vocabolario di glosse operano soprattutto su token
separati da spazi: confondere i due livelli produce gate e vincoli errati.

\section{Supervisione, policy optimization e vincoli}
SFT apprende dai target gold; GRPO confronta più completion dello stesso prompt;
il constraint modifica il supporto disponibile in generazione. Sono interventi
diversi: un Trie può garantire compatibilità lessicale senza migliorare il
contenuto, e un reward ben formulato può essere empiricamente poco informativo.
BLEU \cite{papineni2002bleu}, ROUGE \cite{lin2004rouge} e chrF misurano aspetti
diversi e vanno letti insieme.

\section{Riproducibilità}
Un numero senza identità di dataset, modello, tokenizer, vocabolario, retrieval,
grammar, decoding, seed e budget non è un risultato riproducibile. Il progetto
materializza queste identità negli artifact e usa ID stabili per cache e
bootstrap appaiato. Provenienza e failure espliciti sono parte del metodo.

\chapter{Dati, split e prompting}\label{chap:data}
\section{Preprocessing ASLG-PC12}
Il loader in \pathcode{src/datasets/aslg_dataset.py} legge \pathcode{text} e
\pathcode{gloss}, scarta righe vuote e deduplica prima dello split usando il testo
inglese normalizzato (minuscole e spazi compressi). Dal train grezzo deriva uno
split deterministico 90/10 con seed 42; un test grezzo eventuale non viene usato.
Ogni riga riceve un ID SHA-256 del testo normalizzato e gold, e difficoltà per
lunghezza: 1--5, 6--15, almeno 16 glosse.

\section{Vocabolario e anti-leakage}
\pathcode{data/gloss_vocab.txt} è costruito soltanto dal train, ordinato in modo
deterministico e accompagnato da sidecar; include BOS, EOS e UNK. Dev/test non
definiscono lessico, archi o pesi. La matrice bigram è opzionale e diagnostica.
La stessa separazione vale per l'indice retrieval.

\section{Prompt e retrieval TF--IDF}
\pathcode{src/utils/prompting.py} usa il chat template del tokenizer, con fallback
ChatML. Zero-shot contiene istruzione e sorgente; few-shot antepone esempi
\texttt{English:/ASL gloss:}. Il retriever in
\pathcode{src/retrieval/example_retriever.py} usa TF--IDF, $k=3$, soglia di
autosimilarità 0.98 e cache \pathcode{data/retriever_index}. Le dimostrazioni
provengono solo dal train e il candidato corrente va escluso: controllare ID e
hash evita self-retrieval e contaminazione. Nell'evaluator il nome CLI è
\pathcode{retrieval}; negli artifact è \pathcode{few-shot}.

\chapter{Modello, QLoRA e SFT}\label{chap:sft}
\section{Configurazione esatta}
La base usa Qwen2.5-0.5B-Instruct, sequenza massima 1024, pesi 4-bit, calcolo
bfloat16, Unsloth e una GPU. LoRA modifica
\begin{equation}W=W_0+\frac{\alpha}{r}BA,\end{equation}
con $r=32$, $\alpha=64$, dropout 0.05, seed 3407 e target elencati in
\pathcode{base.yaml}: proiezioni q/k/v/o e gate/up/down \cite{hu2021lora}. Il
backbone quantizzato resta congelato; gli adapter non sono per questo ``4-bit''.

\section{Obiettivo e ricetta SFT}
\begin{equation}\mathcal L_{\rm SFT}=-\sum_{t=1}^{T}\log p_\theta(y_t\mid x,y_{<t}).\end{equation}
La config canonica \pathcode{sft/zero-shot.yaml} usa 3 epoche, batch/device 4,
accumulo 4, learning rate $2\,10^{-5}$, cosine scheduler, warmup 50, eval fraction
0.02 ed early stopping patience 3. SFT è addestrato zero-shot ma valutato con
entrambi i prompt.

\section{Fingerprint e riuso}
Il checkpoint SFT è riusabile soltanto se fingerprint di config, modello,
tokenizer, dataset/split e adapter coincidono. \method{sft-grpo} riusa l'adapter
SFT; non va selezionato un checkpoint solo perché è il più recente. Conservare
config risolta e origine, e usare \pathcode{--resume} solo con la stessa identità.

\chapter{GRPO e Dr-GRPO}\label{chap:grpo}
Per ogni prompt vengono generate $G=8$ completion. Concettualmente il confronto
di gruppo centra i reward, $A_i=r_i-\bar r$. La formula spesso mostrata,
$A_i=(r_i-\bar r)/(s_r+\epsilon)$, \emph{non} descrive questa config: in TRL
0.24.0 \pathcode{scale_rewards: none} disattiva la divisione per deviazione
standard, non il confronto nel gruppo \cite{trlgrpo,shao2024deepseekmath}.
Gruppi a reward costante non danno discriminazione.

Con rapporto token $\rho_{it}=\pi_\theta/\pi_{\rm old}$ e $\epsilon=0.2$, il
surrogate usa
\begin{equation}\min(\rho_{it}A_i,\operatorname{clip}(\rho_{it},1-\epsilon,1+\epsilon)A_i).
\end{equation}
Il KL verso la reference ha coefficiente $\beta=0.04$; clipping rispetto alla
policy dei rollout e ancoraggio alla reference hanno ruoli distinti.

\section{Dr-GRPO e parametri correnti}
\pathcode{loss_type: dr_grpo} usa un denominatore globale costante invece della
media per lunghezza della singola completion, schematicamente
\begin{equation}\mathcal L=-\frac{1}{G L_{\max}}\sum_i\sum_{t\le |y_i|}\mathcal S_{it}.
\end{equation}
Questo riduce un bias, non garantisce lunghezze corrette \cite{liu2025drgrpo}.
Config: batch/device 1, accumulo 8, AdamW paged 8-bit, weight decay 0.1, grad norm
0.1, cosine, max prompt 256, $L_{\max}=128$, temperatura 0.7, importance sampling
token-level e completion troncate mascherate. \method{hot} cambia solo $T$ a 1.3.

\chapter{Reward: definizione, qualificazione e limiti}\label{chap:rewards}
\section{Contratto comune e gate}
Esiste sempre \emph{un solo} reward selezionato da \pathcode{reward.name}, mai una
somma pesata. La completion viene normalizzata con
\pathcode{extract_gloss_text}, poi Unicode \pathcode{casefold} e
\pathcode{split}; la reference è testo piano e riceve casefold+split. Il gate
richiede completion e reference non vuote e ogni token generato nel vocabolario.
Un fallimento restituisce \pathcode{invalid_score}, default $-1$. Il builder lo
accetta soltanto finito e in $[-1,1]$: NaN, infinito e valori esterni sono
rifiutati. La
similarità viene clampata in $[0,1]$ e $R=2s-1$, quindi ogni reward è
rigorosamente in $[-1,1]$.

\section{Le cinque candidate esatte}
\begin{description}[style=nextline]
\item[edit-validity] $s=1-d_{Lev}(g,y)/\max(|g|,|y|)$, Levenshtein sui token.
\item[token-f1-validity] F1 dell'intersezione multinsieme con conteggi clipped;
premia contenuto ma è intenzionalmente invariante al riordino.
\item[chrfpp-validity] SacreBLEU chrF++: char order 6, word order 2, $\beta=2$,
whitespace disattivato e nessun epsilon smoothing.
\item[rouge-l-validity] F1 dalla LCS sui token, senza stemming.
\item[sbleu2-exp-validity] sentence BLEU-2, effective order, smoothing
esponenziale e tokenizer \pathcode{none} \cite{post2018sacrebleu}.
\end{description}
Correttezza della formula non equivale a utilità empirica. Edit penalizza ogni
operazione uniformemente; multiset F1 ignora ordine; chrF++ tollera vicinanza di
caratteri; ROUGE-L conserva ordine ma ammette sottosequenze; BLEU-2 è fragile su
frasi corte, specialmente a una glossa.

\section{Qualificazione congelata}
Prima di un'ablazione si misurano tutte le candidate sullo stesso artifact con
$G=8$, retrieval, Trie, $T=0.7$, massimo 128. Le soglie per candidata sono:
gruppi a varianza zero $\le0.50$, all-min $\le0.50$, score unici medi $\ge1.50$,
best-minus-mean $\ge0.05$, oltre a ranking, lunghezza e perturbazioni. Il report
registra hash input/vocabolario, identità decoding, protocollo e reward, firme
SacreBLEU e digest SHA-256 congiunto di \pathcode{src/rewards/t2g_rewards.py} e
\pathcode{src/utils/text_utils.py}.

Le perturbazioni sono metric-specifiche: l'inversione per token-F1 è invarianza
attesa e non fa gate. No-op, palindromi e casi non valutabili sono contati in
\pathcode{no_op_count} e \pathcode{structurally_unevaluable_count}, non
mascherati come successi o failure e non fanno gate. Degradazioni attese valutabili richiedono
media non negativa e \emph{evidenza positiva} di calo. È screening manuale, non
prova di superiorità.

\section{Hard gate a due livelli}
\pathcode{cluster/train.sh} richiede in \pathcode{EXTRA_ARGS} l'opzione
\pathcode{--reward-qualification-report} per config reward-ablation. In modo
indipendente, il trainer Python valida subito dopo il config e prima di risolvere
run, dati o modello. Rifiuta report mancante/non JSON, protocollo stale,
candidate fallita o ineleggibile, mismatch di digest/impostazioni/SacreBLEU,
gruppo, temperatura, prompt/retrieval, Trie o max length, artifact sorgente
mancante/hash-alterato e vocabolario mancante, stale o non coincidente col
config. L'invocazione Python diretta non aggira il gate; nessun gate promuove
automaticamente una metrica.
\begin{lstlisting}[language=bash]
python -m src.analysis rewards --config experiments/configs/qwen25-05b/probes/rewards.yaml --input experiments/results/.../generations_x.json
EXTRA_ARGS="--reward-qualification-report experiments/analysis/qwen25-05b/rewards/report.json" bash cluster/run_all.sh grpo-few-reward-token-f1
\end{lstlisting}
Le config sono manual-only. TUI/driver mostrano il campo report solo per train
reward, accettano path repository-relative sotto \pathcode{experiments/analysis/}
e nella modalità train+eval inoltrano il report soltanto alla voce train.

\section{Copertura dei test borderline}
La suite copre vocabolario vuoto, gold vuoto o mancante, OOV e casefold Unicode.
Copre inoltre spazi e punteggiatura, operazioni edit, duplicati e riordino,
sequenze disgiunte, SacreBLEU a un token, bound stretti, chat TRL malformata o
con final message, errori config e provenienza. Questi test provano invarianti e
failure mode, non superiorità empirica.

\chapter{Decodifica vincolata}\label{chap:constraints}
\section{Trie primario e maschera del vocabolario}
\warningbox{Non attribuire al Trie il filtraggio di voci numeriche o subtoken. È
\pathcode{GlossVocabularyMask} a filtrare token ID secondo le proprie regole. Il
Trie viene poi compilato sulle stringhe in \pathcode{mask.vocab}, usando le
codifiche normali e con spazio prefisso. Sono due stadi distinti.}
Con storia $h_t$ e insieme ammesso $A(h_t)$, il logits processor pone a
$-\infty$ i logit fuori insieme e rinormalizza. Al termine di una glossa consente
la successiva o EOS e gestisce separatori/newline. Il vincolo garantisce un
percorso compilato, non correttezza semantica.

\section{PDA manuale}
\pathcode{ablations/sft-grpo-zero-pda.yaml} abilita GrammarLLM PDA e lookahead.
La grammatica LL(1) rappresenta prima glossa, coda e EOS; lo stack di parsing è
diverso dal cammino prefisso del Trie. Il confronto mantiene fissi modello,
inizializzazione, dati, prompt, seed e budget e richiede il gate sull'intero
vocabolario reale. PDA non è il percorso di produzione né un quinto metodo.

\section{Massa mascherata}
Prima del masking,
\begin{equation}M_t=\sum_{v\in A(h_t)}\operatorname{softmax}(z_t)_v,\qquad
R_t=1-M_t.\end{equation}
Massa ammessa/rimossa ed entropia condizionata descrivono il conflitto tra policy
e vincolo. Sono diagnostici congelati, non reward né obiettivi trainabili. Valori
bassi di $M_t$ non provano che il constraint corregga il significato.

\chapter{Benchmark structured-loss a stati ridotti}\label{chap:structured}
È un benchmark isolato: congela Qwen, estrae il vettore dell'ultimo layer al
confine assistant e allena soltanto una testa di emissione posizionale. Non è
collegato al trainer SFT/GRPO, al driver o alla TUI.

Lo spazio ha le 512 glosse train più frequenti più \pathcode{<OTHER>}; BOS/EOS
esistono solo nel grafo. Sequenze oltre 64 sono escluse e registrate, mai
troncate. Il grafo usa $\alpha=0.1$ e scala 0.25; dev/test non definiscono stati o
archi. Dichiarando $y_0=\mathrm{BOS}$, lo score comprende l'arco EOS finale:
\begin{equation}S(y;x,L)=\sum_{t=1}^{L}e_t(y_t;x)+
\sum_{t=1}^{L}a(y_{t-1},y_t)+a(y_L,\mathrm{EOS}).\end{equation}
Con $\log Z=\log\sum_{|z|=L}\exp S(z;x,L)$, l'implementazione riporta
\begin{equation}\mathrm{NLL}=(\log Z-S)/L,\end{equation}
non la somma non normalizzata. La partizione termina esattamente a $L$.

Le braccia sono CE indipendente, uniform-support, true-weights e
shuffled-weights a supporto fisso. Feature e righe sono condivise; argmax
posizionale a lunghezza gold non è Viterbi. Le metriche di estrazione stanno nel
manifest separato dalle metriche delle teste. La resource gate segnala pass/fail
ma non promuove automaticamente: true-weights deve battere i controlli su tutti
i seed configurati e la decisione resta manuale. Avvio offline:
\begin{lstlisting}[language=bash]
sbatch cluster/structured_probe.sh
\end{lstlisting}

\chapter{Valutazione e inferenza}\label{chap:evaluation}
\section{Deployment e sampling}
Su 2.000 esempi test seeded, ogni prompt ha una completion \method{deployment}
greedy (\pathcode{do_sample=false}) e $N=5$ draw \method{sampling} indipendenti a
$T=0.7$. Le headline sono solo greedy; sampling misura qualità attesa,
diversità/headroom e Pass@k. Un oracle scelto col gold è non deployable e va
separato. I seed per sample derivano deterministicamente da seed di protocollo,
ID stabile sorgente+gold, mode e indice: batching, cache e resume non li cambiano.

\section{Metriche e interpretazione}
Si riportano corpus BLEU-4 e chrF2, ROUGE-L medio, micro token-F1, exact match
normalizzato, edit similarity, validità e lunghezze. Il successo Pass@k richiede
ROUGE-L $\ge0.3$ e validità lessicale; lo stimatore è
$1-\binom{n-c}{k}/\binom nk$ \cite{chen2021codex}. Includere errore di lunghezza
signed/assoluto, ratio, empty e truncation-hit. Validità non è qualità; una
decodifica vincolata può aumentare la prima e peggiorare le altre metriche.

\section{Intervalli e confronti}
Il bootstrap è deterministico e clustered per prompt. Le metriche sentence dei
draw sono prima aggregate nel prompt; BLEU/chrF corpus vengono ricalcolati a ogni
resample. Un delta appaiato richiede gli stessi ID unici nello stesso ordine e
gli stessi indici bootstrap. Non usare intervalli indipendenti per sostituire il
test appaiato.

\section{Identità della cache}
Cache e baseline esplicita devono coincidere su versione metrica, hash degli ID,
dataset, modello, tokenizer, vocabolario, retrieval, grammar, decoding, seed,
mode, $N$, budget e temperatura. Ogni mismatch fallisce esplicitamente. Gli
artifact conservano \pathcode{decoding_mode}, \pathcode{completion_index}, raw
generations, metriche, config, checkpoint e firme SacreBLEU.

\chapter{Disegno sperimentale canonico}\label{chap:experiments}
\section{Sette celle e dodici voci}
Le sette celle sono due baseline eval-only (base zero/few), SFT zero-shot e il
$2\times2$ GRPO: inizializzazione base/SFT per prompt train zero/few. La coda ha
12 voci: 2 eval baseline più train+eval per le 5 celle apprese. Ogni eval appresa
esegue zero-shot e retrieval. Non contare le due gambe come metodi diversi.

\begin{center}\begin{tabular}{lll}\toprule
Metodo & Inizializzazione & Prompt train\\\midrule
SFT & base & zero-shot\\
GRPO & base & zero-shot\\
GRPO & base & few-shot\\
SFT-GRPO & SFT & zero-shot\\
SFT-GRPO & SFT & few-shot\\\bottomrule
\end{tabular}\end{center}

\section{Config e artifact}
La radice è \pathcode{experiments/configs/qwen25-05b/}; \pathcode{base.yaml} è
ricetta condivisa, non cella. I runnable dichiarano \pathcode{model_tag},
\pathcode{method}, \pathcode{train_prompt_mode}, \pathcode{variant} e
\pathcode{kind}. Checkpoint/log seguono
\begin{lstlisting}
experiments/{checkpoints,logs}/qwen25-05b/<method>/<train-prompt>/[ablations/<variant>/]run_<timestamp>/
\end{lstlisting}
e risultati
\begin{lstlisting}
experiments/results/qwen25-05b/<method>/<train-prompt>/[ablations/<variant>/]eval-<prompt>/run_<timestamp>/
\end{lstlisting}
Le baseline usano \pathcode{baseline/<prompt>/run_*}. Identità config e artifact
devono concordare; non copiare YAML in gerarchie parallele.

\section{Ablazioni fuori matrice}
PDA, hot e cinque reward sono manuali. Rollout/reward/Markov sono probe su
generazioni congelate. Structured-loss è un benchmark head-only. Nessuno entra
nella campagna primaria senza nuova preregistrazione.

\chapter{Diagnostica congelata e decisioni}\label{chap:diagnostics}
\section{Rollout e reward}
I probe consumano \pathcode{generations_*.json}; non cercano input e non generano
automaticamente. Rollout raggruppa per ID, controlla gold coerente e misura
varianza, saturazione, unicità, best-minus-mean, correlazioni e bucket di
lunghezza. Le boundary del reward sono quelle del \cref{chap:rewards}; superarle
autorizza solo l'ispezione/ablazione manuale.

\section{Markov, Viterbi hard e soft}
Il bigram $O(L)$ può usare il lessico completo; hard/soft usano piccoli grafi e
falliscono prima dell'allocazione oltre \pathcode{max_states=256}. Il Viterbi hard
massimizza transizioni logaritmiche e applica $-\lambda\mathbf1[i=j]$ soltanto
alle transizioni \emph{interiori}; l'arco finale verso EOS non riceve la penalità
di self-loop. Soft sostituisce max con log-sum-exp temperata. Si analizzano
Spearman, discriminazione pairwise, lunghezza e correlazione parziale.
In formula, solo per i passi interni,
\begin{equation}
\delta_t(j)=\max_i\{\delta_{t-1}(i)+\log P_{ij}-\lambda\mathbf1[i=j]\};
\end{equation}
la chiusura usa $\max_i\{\delta(i)+\log P_{i,\mathrm{EOS}}\}$ senza penalità.

Non è una Differentiable Viterbi Layer: non vi sono emissioni apprese,
backpointer differenziabili o gradienti \cite{viterbiplannet2026}. La decisione è
se il segnale aggiunga ranking oltre edit/bigram sui medesimi sample; anche in
caso positivo serve un'ablazione di training isolata per una pretesa causale.
\begin{lstlisting}[language=bash]
INPUT=experiments/results/.../generations_x.json sbatch cluster/probe.sh rollouts experiments/configs/qwen25-05b/probes/rollouts.yaml
INPUT=experiments/results/.../generations_x.json sbatch cluster/probe.sh markov experiments/configs/qwen25-05b/probes/markov.yaml
\end{lstlisting}

\chapter{Guida alla codebase ed estensioni}\label{chap:codebase}
\section{Mappa per directory e simboli}
\begin{description}[style=nextline]
\item[\pathcode{src/datasets/}] prepara righe, split, matrici e grafi strutturati.
\item[\pathcode{src/models/}] \pathcode{model_loader.py} carica Qwen/adapter;
\pathcode{structured_gloss_head.py} implementa la testa ridotta.
\item[\pathcode{src/training/}] entry point SFT, GRPO, eval e benchmark strutturato.
\item[\pathcode{src/grammar/}] implementa logits processor Trie/PDA e telemetria.
\item[\pathcode{src/rewards/t2g_rewards.py}] registry, normalizzazione, gate e metriche reward.
\item[\pathcode{src/retrieval/}] indice e ricerca TF--IDF.
\item[\pathcode{src/analysis/}] rollout, reward, Markov e structured benchmark.
\item[\pathcode{src/utils/}] config, path/identity, prompt, metriche, cache e monitor.
\item[\pathcode{experiments/configs/}] ricette canoniche; \pathcode{cluster/} launcher
offline/chain; \pathcode{remote/} API e TUI; \pathcode{tests/} invarianti.
\end{description}

\section{Ricette di estensione}
\paragraph{Nuova config.} Estendere la base, cambiare un asse, assegnare identità
univoca, aggiungere validazione/path test e registrarla nel driver solo se deve
essere esposta. Non confondere prompt e method.
\paragraph{Nuovo reward.} Implementare similarità pura, riusare normalizzazione,
gate, clamp e transform, registrare un solo nome, aggiungere borderline e
perturbation contract, config manuale e qualificazione con digest/hash.
\paragraph{Nuovo probe.} Consumare artifact congelato esplicito, versionare
schema/protocollo, separare diagnostica da training e definire boundary prima di
guardare i risultati.
\paragraph{Nuovo backend retrieval.} Implementare build/query, fingerprint di
corpus e parametri, esclusione self-match, cache fail-closed e test anti-leakage.
\paragraph{Nuova route/API.} Aggiungere modello request, autenticazione,
validazione repository-relative, operazione idempotente e test API/TUI; non
considerare SQLite fonte autoritativa.

\section{Validazione}
La suite comprende config, dataset, prompt/retrieval, reward, grammar, metriche,
training/eval, cache/path, launcher offline, probe, API e TUI. Comandi canonici:
\begin{lstlisting}[language=bash]
uv run python -m pytest tests/ -v
python tests/validate_configs.py
\end{lstlisting}
Uno skip GPU/cache è un prerequisito assente, non un pass. Registrare comando,
versioni, pass/fail/skip e non riportare conteggi storici come correnti.

\chapter{Troubleshooting, limiti e passi successivi}\label{chap:troubleshooting}
\section{Diagnosi per sintomo}
\begin{longtable}{p{.28\textwidth}p{.60\textwidth}}\toprule
Sintomo & Azione\\\midrule
Config/identity mismatch & Stampare config risolta; verificare method, prompt,
variant, kind e checkpoint source; non rinominare directory a mano.\\
Cache/model offline mancante & Non abilitare rete nel job: rieseguire setup sul
login node, quindi preflight.\\
Output vuoti o tutti invalidi & Ispezionare chat extraction, EOS, vocabolario,
OOV e massa ammessa; distinguere gate reward da Trie.\\
Reward senza varianza & Controllare gruppi, gold, $G=8$, gate e distribuzione;
non lanciare training per ``vedere se migliora''.\\
OOM & Verificare config effettiva, batch/accumulo, max length, quantizzazione e
processi concorrenti; non cambiare più assi nell'esperimento.\\
Eval non appaiabile & Rigenerare con stessi ID ordinati e identica protocol
identity; non forzare la cache.\\
PDA fallisce & Eseguire preflight con \pathcode{PDA=1} sul vocabolario completo.\\
W\&B non visibile & Il compute è offline: preservare directory e sincronizzare
dopo, da nodo con rete.\\\bottomrule
\end{longtable}

\section{LEGACY --- NON CONFRONTABILE}
\warningbox{Artifact, metriche o figure anteriori al protocollo canonico
\pathcode{qwen25-05b}, incluse versioni metriche precedenti, sono LEGACY E NON
CONFRONTABILI. Non unirli a tabelle correnti.}
Questo manuale omette valori storici: non sono evidenza della campagna attuale.

\section{Passi successivi responsabili}
Direzioni possibili, non implementazioni presenti: valutazione umana/linguistica
ASL; più seed e potenza; retrieval appreso con anti-leakage; reward d'ordine
qualificato e isolato; transizioni condizionate alla sorgente; vera testa
strutturata con Viterbi e gradient-flow test; studio causale della massa rimossa.
Ogni estensione richiede ipotesi, controllo, failure mode, test, budget e criterio
di decisione preregistrati. Superare un probe non dimostra beneficio di training.

\chapter{Operazioni sul cluster}\label{chap:cluster}
\section{Setup online sul login node}
\begin{lstlisting}[language=bash]
ssh <user>@gcluster.dmi.unict.it
cd ~/neuro_symbolic_t2g
bash cluster/setup.sh
\end{lstlisting}
Il setup online esegue Apptainer \pathcode{/shared/sifs/latest.sif} direttamente
sul login node: niente \pathcode{srun}, allocazione GPU o Python host. Installa in
\pathcode{~/.local} senza alterare lo stack torch/CUDA, verifica versioni, scarica
Qwen e ASLG-PC12 e costruisce vocabolario/sidecar. \pathcode{BUILD_BIGRAM=1}
abilita la matrice opzionale.
Le versioni presidiate includono Transformers 5.3.0, TRL 0.24.0, PEFT 0.19.1,
Unsloth 2026.7.1 e torchao 0.16--0.17.

\section{Compute rigorosamente offline}
Preflight, train, eval e probe esportano prima di Python:
\begin{lstlisting}
HF_HUB_OFFLINE=1
TRANSFORMERS_OFFLINE=1
HF_DATASETS_OFFLINE=1
WANDB_MODE=offline
WANDB_DISABLE_WEAVE=true
WANDB_SILENT=true
\end{lstlisting}
Il preflight controlla import, snapshot, dataset, vocab, retrieval/bigram se
richiesti e W\&B. Non scarica e non allena.

\section{Coda e campaign chain}
La QoS ammette un solo job submitted/attivo per utente: \pathcode{run_all.sh}
mantiene una chain sequenziale di 12 voci. Non lanciare una seconda campagna in
parallelo. Comandi:
\begin{lstlisting}[language=bash]
CONFIG=experiments/configs/qwen25-05b/sft-grpo/zero-shot.yaml sbatch cluster/preflight.sh
bash cluster/run_all.sh
CONFIG=experiments/configs/qwen25-05b/sft-grpo/zero-shot.yaml EXTRA_ARGS="--resume" sbatch cluster/train.sh
\end{lstlisting}
Usare gli strumenti \pathcode{chain-show}/monitor per osservare l'avanzamento.
Una failure va corretta e ripresa con la stessa identità, non aggirata cambiando
path. PDA e reward usano i gate aggiuntivi descritti nei capitoli 6 e 7.

\chapter{API remota, TUI e monitoraggio}\label{chap:remote}
Il servizio FastAPI in \pathcode{remote/} esegue via SSH tick idempotenti. Lo
stato autoritativo è \pathcode{.chain_state/} sul cluster; SQLite remoto è solo
cache/journal. Pause è soft e non cancella il job attivo. Tutte le route tranne
\pathcode{GET /} richiedono \pathcode{X-Auth-Token}.

\begin{lstlisting}[language=bash]
uv run --extra dev uvicorn app:app --port 8000 --app-dir remote
uv run --extra tui python remote/tui.py --url http://localhost:8000 --token local-token
\end{lstlisting}
Il registry espone celle canoniche e varianti manuali. Il campo storico API
\pathcode{ablation:true} seleziona/sostituisce la coda predefinita, non include
automaticamente PDA/hot. Per reward train il report qualificato è obbligatorio;
eval-only non deve riceverlo.
Rollout e Markov non sono queue config; il benchmark structured è assente da
driver e TUI. Un 401 indica autenticazione, mentre un 502 segnala tipicamente un
problema SSH/cluster; lo status può restituire cache marcata non raggiungibile.

\section{W\&B offline e sicurezza}
I compute scrivono run W\&B offline in directory isolate per gamba eval. Il sync
è successivo da ambiente con rete \cite{wandboffline}. Non mettere token HF,
W\&B, API o chiavi SSH in repository, YAML, log o command history condivisa;
preferire environment/secret store, permessi minimi e rotazione. Esporre l'API
solo su interfaccia fidata o tunnel SSH, usare token forte e non stampare header.
Validare ogni path come repository-relative e non accettare argomenti shell
arbitrari dalla TUI.

Da Windows gli script \pathcode{sync_cluster.ps1} scaricano log, checkpoint e
risultati. Conservare la gerarchia timestampata: appiattirla perde provenienza.

\appendix
\chapter{Comandi canonici e glossario}\label{app:commands}
\section{Comandi}
Usare i path completi nei file di job; i segnaposto vanno sostituiti senza
cambiare l'identità sperimentale.
\begin{lstlisting}[language=bash]
# login node online, dentro Apptainer gestito dallo script
bash cluster/setup.sh

# compute offline
CONFIG=experiments/configs/qwen25-05b/sft-grpo/zero-shot.yaml sbatch cluster/preflight.sh
bash cluster/run_all.sh
CONFIG=experiments/configs/qwen25-05b/sft/zero-shot.yaml sbatch cluster/train.sh
CONFIG=experiments/configs/qwen25-05b/baseline/few-shot.yaml sbatch cluster/eval.sh

# manuali
PDA=1 CONFIG=experiments/configs/qwen25-05b/ablations/sft-grpo-zero-pda.yaml sbatch cluster/preflight.sh
bash cluster/run_all.sh sft-grpo-zero-pda
bash cluster/run_all.sh sft-grpo-zero-hot

# validazione locale
uv run python -m pytest tests/ -v
python tests/validate_configs.py

# documento
latexmk -pdf main.tex
\end{lstlisting}

\section{Glossario}
\begin{description}[style=nextline]
\item[Artifact identity] Config, modello/checkpoint, dati, tokenizer, prompt,
constraint, seed, budget, hash e versioni che rendono interpretabile un output.
\item[ASL gloss] Etichetta testuale semplificata di un segno, non lingua completa.
\item[ASLG-PC12] Corpus parallelo del progetto.
\item[Deployment] Una completion greedy per prompt; sorgente delle headline.
\item[Diagnostico congelato] Analisi senza gradiente su generazioni esistenti.
\item[Dr-GRPO] Loss con denominatore globale per ridurre length bias.
\item[Few-shot/retrieval] Modalità di prompt con esempi TF--IDF, non metodo.
\item[GRPO] Policy optimization relativa a gruppi di completion.
\item[LoRA/QLoRA] Adapter low-rank su backbone congelato/quantizzato.
\item[Metodo] Base, SFT, GRPO o SFT-GRPO.
\item[PDA] Automa a pila GrammarLLM, ablazione manuale.
\item[Prompt-clustered bootstrap] Resampling dei prompt mantenendo uniti i draw.
\item[Reward] Singolo segnale scalare usato per aggiornare la policy.
\item[Sampling] Cinque draw eval a $T=0.7$; la qualificazione reward ne richiede otto.
\item[SFT] Fine-tuning supervisionato sui target gold.
\item[Structured NLL] NLL con partizione a lunghezza gold fissa; qui $(\log Z-S)/L$.
\item[Trie] Albero di prefissi compilato dal vocabolario mascherato; percorso primario.
\item[Validità] Appartenenza lessicale/formale, distinta dalla qualità traduttiva.
\item[Viterbi] DP hard/soft qui diagnostica a piccoli stati, non DVL addestrabile.
\item[Zero-shot] Modalità di prompt senza esempi, non metodo.
\end{description}

\backmatter
\bibliographystyle{plain}
\bibliography{references}
\end{document}
